\documentclass[11pt,letterpaper]{article}

% ── Packages ──────────────────────────────────────────────────
\usepackage[margin=1in]{geometry}
\usepackage{amsmath,amssymb,amsfonts}
\usepackage{algorithm}
\usepackage{algpseudocode}
\usepackage{booktabs}
\usepackage{graphicx}
\usepackage{xcolor}
\usepackage{hyperref}
\usepackage{listings}
\usepackage{fancyhdr}
\usepackage{enumitem}
\usepackage{tikz}
\usepackage{tcolorbox}
\usepackage{caption}

\tcbuselibrary{listings,skins,breakable}

% ── Colors ────────────────────────────────────────────────────
\definecolor{zetablue}{RGB}{30,70,130}
\definecolor{codebg}{RGB}{248,248,248}
\definecolor{codeborder}{RGB}{200,200,200}
\definecolor{keywordblue}{RGB}{0,0,180}
\definecolor{commentgreen}{RGB}{0,128,0}
\definecolor{stringred}{RGB}{163,21,21}
\definecolor{warnred}{RGB}{180,30,30}
\definecolor{notebg}{RGB}{240,248,255}
\definecolor{warnbg}{RGB}{255,245,238}

% ── Hyperref ──────────────────────────────────────────────────
\hypersetup{
    colorlinks=true,
    linkcolor=zetablue,
    citecolor=zetablue,
    urlcolor=zetablue,
}

% ── Code listings ─────────────────────────────────────────────
\lstdefinestyle{pythonstyle}{
    language=Python,
    basicstyle=\ttfamily\small,
    keywordstyle=\color{keywordblue}\bfseries,
    commentstyle=\color{commentgreen}\itshape,
    stringstyle=\color{stringred},
    backgroundcolor=\color{codebg},
    frame=single,
    rulecolor=\color{codeborder},
    numbers=left,
    numberstyle=\tiny\color{gray},
    numbersep=8pt,
    tabsize=4,
    showstringspaces=false,
    breaklines=true,
    xleftmargin=2em,
    framexleftmargin=1.5em,
}

\lstdefinestyle{cudastyle}{
    language=C++,
    basicstyle=\ttfamily\small,
    keywordstyle=\color{keywordblue}\bfseries,
    commentstyle=\color{commentgreen}\itshape,
    stringstyle=\color{stringred},
    backgroundcolor=\color{codebg},
    frame=single,
    rulecolor=\color{codeborder},
    numbers=left,
    numberstyle=\tiny\color{gray},
    numbersep=8pt,
    tabsize=4,
    showstringspaces=false,
    breaklines=true,
    xleftmargin=2em,
    framexleftmargin=1.5em,
    morekeywords={__global__,__device__,__shared__,__constant__,__ldg,
                  dim3,cudaMalloc,cudaFree,cudaMemcpy,cudaMemcpyToSymbol,
                  cudaDeviceSynchronize,cudaEventCreate,cudaEventRecord,
                  cudaEventSynchronize,cudaEventElapsedTime,blockIdx,
                  blockDim,threadIdx,gridDim,float2,float4,int2},
}

% ── Header/Footer ─────────────────────────────────────────────
\pagestyle{fancy}
\fancyhf{}
\fancyhead[L]{\small\textcolor{zetablue}{\textbf{Zetascope}} --- GPU Engineer Interview}
\fancyhead[R]{\small\textcolor{gray}{Confidential}}
\fancyfoot[C]{\small\thepage}
\renewcommand{\headrulewidth}{0.4pt}
\renewcommand{\footrulewidth}{0pt}

% ── Custom boxes ──────────────────────────────────────────────
\newtcolorbox{notebox}[1][]{
    colback=notebg,
    colframe=zetablue,
    fonttitle=\bfseries,
    title={#1},
    arc=2pt,
    boxrule=0.5pt,
    left=8pt, right=8pt, top=4pt, bottom=4pt,
}

\newtcolorbox{warnbox}[1][]{
    colback=warnbg,
    colframe=warnred,
    fonttitle=\bfseries\color{warnred},
    title={#1},
    arc=2pt,
    boxrule=0.5pt,
    left=8pt, right=8pt, top=4pt, bottom=4pt,
}

% ── Title ─────────────────────────────────────────────────────
\title{%
    \vspace{-1cm}
    {\Large\textcolor{zetablue}{\textbf{Zetascope, Inc.}}}\\[0.3cm]
    {\LARGE GPU Engineer Technical Interview}\\[0.2cm]
    {\Large Brute-Force Track-Before-Detect on CUDA}\\[0.3cm]
    \rule{\textwidth}{0.4pt}
}
\author{}
\date{\today}

% ==============================================================
\begin{document}
\maketitle
\thispagestyle{fancy}

% ──────────────────────────────────────────────────────────────
\section{Introduction}
% ──────────────────────────────────────────────────────────────

\subsection{Background}

Track-Before-Detect (TBD) is a signal processing paradigm for detecting
objects that are too faint to observe in any single measurement. Instead
of detecting first and tracking second (the classical approach), TBD
accumulates evidence across many measurements along candidate motion
hypotheses, exploiting the fact that signal grows as $\mathcal{O}(N)$
while noise grows as $\mathcal{O}(\sqrt{N})$ for $N$ integrated frames.

In the context of \textbf{space domain awareness}, TBD enables detection
of extremely faint resident space objects (RSOs) at cislunar distances,
where the object may contribute only a handful of photons per frame
against a noisy background.

\subsection{The Shift-and-Stack Algorithm}

The simplest TBD method is \textbf{brute-force shift-and-stack}: for every
candidate trajectory defined by a starting pixel position $(x_0, y_0)$ and
constant velocity $(v_x, v_y)$, sum the pixel values along the predicted
path across all frames. The trajectory producing the highest sum (the
\emph{test statistic}) is the detection.

Formally, given a sequence of $N$ frames $\{I_t\}_{t=0}^{N-1}$, the
test statistic is:
\begin{equation}
    \boxed{
    T(x_0, y_0, v_x, v_y)
    \;=\;
    \sum_{t=0}^{N-1} I_t\!\left[\,
        \mathrm{round}(y_0 + v_y \cdot t),\;
        \mathrm{round}(x_0 + v_x \cdot t)
    \,\right]
    }
    \label{eq:test_stat}
\end{equation}

\noindent where $I_t[r, c]$ denotes the pixel value at row $r$, column $c$
of frame $t$. The detection is:
\begin{equation}
    (\hat{x}_0,\, \hat{y}_0,\, \hat{v}_x,\, \hat{v}_y)
    \;=\;
    \arg\max_{(x_0, y_0, v_x, v_y)} \; T(x_0, y_0, v_x, v_y)
    \label{eq:detection}
\end{equation}

\begin{notebox}[Why Brute-Force?]
While more sophisticated algorithms (temporal-velocity pyramids, FFT-based
methods, multi-resolution search) exist to accelerate TBD, the brute-force
approach is the \emph{gold standard} baseline: it is exact, trivially
parallelizable, and serves as the correctness reference for all optimized
methods. Understanding how to efficiently map this embarrassingly parallel
problem onto GPU hardware is a fundamental skill.
\end{notebox}

% ──────────────────────────────────────────────────────────────
\section{Problem Specification}
\label{sec:problem}
% ──────────────────────────────────────────────────────────────

\subsection{Synthetic Data Model}

Each frame is generated as independent Poisson noise with a faint moving
point source injected along a ground-truth trajectory:
\begin{equation}
    I_t[r, c] \;\sim\;
    \mathrm{Poisson}\!\left(\lambda_{\mathrm{bg}}\right)
    + \delta_{r, r_t}\,\delta_{c, c_t}\,
      \mathrm{Poisson}\!\left(\lambda_{\mathrm{obj}}\right)
    \label{eq:data_model}
\end{equation}
where $(c_t, r_t) = \bigl(\mathrm{round}(x_0^* + v_x^* t),\;
\mathrm{round}(y_0^* + v_y^* t)\bigr)$ is the true object position at frame $t$.

\subsection{Problem Parameters}

Table~\ref{tab:params} lists all fixed parameters for this exercise.

\begin{table}[h]
\centering
\caption{Problem parameters.}
\label{tab:params}
\begin{tabular}{@{}llrl@{}}
\toprule
\textbf{Parameter} & \textbf{Symbol} & \textbf{Value} & \textbf{Unit} \\
\midrule
Image width  & $N_x$                 & 256   & pixels \\
Image height & $N_y$                 & 256   & pixels \\
Number of frames & $N$               & 16    & frames \\
Background rate  & $\lambda_{\mathrm{bg}}$ & 100.0 & photons/pixel/frame \\
Object signal    & $\lambda_{\mathrm{obj}}$& 20.0  & photons/frame \\
\midrule
True starting position & $(x_0^*,\, y_0^*)$    & $(128,\, 100)$ & pixels \\
True velocity          & $(v_x^*,\, v_y^*)$    & $(1.5,\, {-}0.75)$ & pixels/frame \\
\midrule
Velocity search min    & $v_{\min}$     & $-3.0$  & pixels/frame \\
Velocity search max    & $v_{\max}$     & $+3.0$  & pixels/frame \\
Velocity step size     & $\Delta v$     & $0.25$  & pixels/frame \\
\midrule
RNG seed               & ---            & 42     & --- \\
\bottomrule
\end{tabular}
\end{table}

\subsection{Signal-to-Noise Analysis}

The per-frame signal-to-noise ratio for a Poisson background is:
\begin{equation}
    \mathrm{SNR}_1
    \;=\; \frac{\lambda_{\mathrm{obj}}}{\sqrt{\lambda_{\mathrm{bg}}}}
    \;=\; \frac{20}{\sqrt{100}}
    \;=\; 2.0
\end{equation}

This is well below typical single-frame detection thresholds
($\mathrm{SNR} \geq 5$). After coherent integration across $N = 16$ frames,
the stacked SNR becomes:
\begin{equation}
    \mathrm{SNR}_N
    \;=\; \mathrm{SNR}_1 \cdot \sqrt{N}
    \;=\; 2.0 \times \sqrt{16}
    \;=\; 8.0
\end{equation}

\noindent which is comfortably above threshold --- but \emph{only if the
correct trajectory is used for integration}. This is the fundamental
premise of TBD.

\subsection{Search Space Complexity}

The velocity grid contains:
\begin{equation}
    N_v \;=\; \left\lfloor \frac{v_{\max} - v_{\min}}{\Delta v} \right\rfloor + 1
    \;=\; \left\lfloor \frac{6.0}{0.25} \right\rfloor + 1
    \;=\; 25
\end{equation}

\noindent velocity steps per axis, giving $N_v^2 = 625$ velocity hypotheses.
The total number of candidate trajectories is:
\begin{equation}
    N_{\mathrm{traj}}
    \;=\; N_x \cdot N_y \cdot N_v^2
    \;=\; 256 \times 256 \times 625
    \;=\; 40{,}960{,}000
    \;\approx\; 41\text{M}
\end{equation}

Each trajectory evaluation requires $N = 16$ pixel lookups, yielding:
\begin{equation}
    N_{\mathrm{reads}}
    \;=\; N_{\mathrm{traj}} \times N
    \;=\; 40{,}960{,}000 \times 16
    \;=\; 655{,}360{,}000
    \;\approx\; 655\text{M pixel reads}
\end{equation}

% ──────────────────────────────────────────────────────────────
\section{Algorithm Pseudocode}
\label{sec:algorithm}
% ──────────────────────────────────────────────────────────────

\begin{algorithm}[H]
\caption{Brute-Force Shift-and-Stack TBD}
\label{alg:tbd}
\begin{algorithmic}[1]
\Require Frames $\{I_t\}_{t=0}^{N-1}$, velocity grid $\mathcal{V} = \{v_{\min}, v_{\min}+\Delta v, \ldots, v_{\max}\}$
\Ensure Detection $(\hat{x}_0, \hat{y}_0, \hat{v}_x, \hat{v}_y)$ and test statistic $\hat{T}$
\State $\hat{T} \gets -\infty$
\For{$y_0 = 0$ \textbf{to} $N_y - 1$}
    \For{$x_0 = 0$ \textbf{to} $N_x - 1$}
        \For{$v_x \in \mathcal{V}$}
            \For{$v_y \in \mathcal{V}$}
                \State $T \gets 0$;\quad $\mathit{valid} \gets \textsc{True}$
                \For{$t = 0$ \textbf{to} $N - 1$}
                    \State $p_x \gets \mathrm{round}(x_0 + v_x \cdot t)$
                    \State $p_y \gets \mathrm{round}(y_0 + v_y \cdot t)$
                    \If{$p_x \notin [0, N_x)$ \textbf{or} $p_y \notin [0, N_y)$}
                        \State $\mathit{valid} \gets \textsc{False}$;\quad \textbf{break}
                    \EndIf
                    \State $T \gets T + I_t[p_y,\, p_x]$
                \EndFor
                \If{$\mathit{valid}$ \textbf{and} $T > \hat{T}$}
                    \State $\hat{T} \gets T$;\quad
                           $(\hat{x}_0, \hat{y}_0, \hat{v}_x, \hat{v}_y) \gets (x_0, y_0, v_x, v_y)$
                \EndIf
            \EndFor
        \EndFor
    \EndFor
\EndFor
\State \Return $(\hat{x}_0, \hat{y}_0, \hat{v}_x, \hat{v}_y, \hat{T})$
\end{algorithmic}
\end{algorithm}

\begin{notebox}[Boundary Handling]
If the predicted position $(p_x, p_y)$ at \emph{any} frame $t$ falls
outside the image boundaries $[0, N_x) \times [0, N_y)$, the entire
trajectory is marked invalid and skipped. This is physically correct:
partial trajectories would have reduced integrated signal and should not
compete with full-length trajectories.
\end{notebox}

% ──────────────────────────────────────────────────────────────
\section{Binary Data Format}
\label{sec:data_format}
% ──────────────────────────────────────────────────────────────

The Python script \texttt{tbd\_problem.py} generates two binary files.
Your CUDA implementation must read these files.

\subsection{\texttt{tbd\_frames.bin}}

\begin{table}[h]
\centering
\caption{Frame data binary layout.}
\begin{tabular}{@{}llll@{}}
\toprule
\textbf{Offset (bytes)} & \textbf{Type} & \textbf{Count} & \textbf{Description} \\
\midrule
0    & \texttt{int32}   & 1   & $N_x$ (image width) \\
4    & \texttt{int32}   & 1   & $N_y$ (image height) \\
8    & \texttt{int32}   & 1   & $N$ (number of frames) \\
12   & \texttt{float32} & $N \times N_y \times N_x$ & Pixel data \\
\bottomrule
\end{tabular}
\end{table}

\noindent The pixel data is stored in \textbf{frame-major, row-major} order:
\begin{equation}
    \texttt{offset}(t, r, c) \;=\; 12 + 4 \times \bigl(t \cdot N_y \cdot N_x + r \cdot N_x + c\bigr) \;\text{bytes}
\end{equation}

\subsection{\texttt{tbd\_params.bin}}

\begin{table}[h]
\centering
\caption{Parameter binary layout.}
\begin{tabular}{@{}llll@{}}
\toprule
\textbf{Offset (bytes)} & \textbf{Type} & \textbf{Count} & \textbf{Description} \\
\midrule
0  & \texttt{float32} & 1 & $v_{\min}$ \\
4  & \texttt{float32} & 1 & $v_{\max}$ \\
8  & \texttt{float32} & 1 & $\Delta v$ \\
\bottomrule
\end{tabular}
\end{table}

% ──────────────────────────────────────────────────────────────
\section{Your Task}
\label{sec:task}
% ──────────────────────────────────────────────────────────────

\subsection{Deliverables}

Implement a \textbf{CUDA/C++ program} that:
\begin{enumerate}[leftmargin=2em]
    \item Reads \texttt{tbd\_frames.bin} and \texttt{tbd\_params.bin}
    \item Performs the brute-force shift-and-stack search
          (Eq.~\ref{eq:test_stat}) entirely on the GPU
    \item Reports the detected trajectory
          $(\hat{x}_0, \hat{y}_0, \hat{v}_x, \hat{v}_y)$ and test
          statistic $\hat{T}$
    \item Reports wall-clock kernel execution time using CUDA events
    \item Compiles with \texttt{nvcc} (any SM architecture $\geq 7.0$)
\end{enumerate}

\subsection{Setup Instructions}

\begin{enumerate}[leftmargin=2em]
    \item Generate the test data:
\begin{lstlisting}[style=pythonstyle, numbers=none, xleftmargin=1em, framexleftmargin=0em]
python3 tbd_problem.py
\end{lstlisting}
    \item Verify that \texttt{tbd\_frames.bin}, \texttt{tbd\_params.bin},
          and \texttt{tbd\_reference.txt} are created.
    \item Implement your solution and compile:
\begin{lstlisting}[style=cudastyle, numbers=none, xleftmargin=1em, framexleftmargin=0em]
nvcc -O3 -arch=sm_86 -o tbd_solution your_solution.cu
./tbd_solution tbd_frames.bin tbd_params.bin
\end{lstlisting}
    \item Your output must match \texttt{tbd\_reference.txt}.
\end{enumerate}

\subsection{Reference Answer}

\begin{table}[h]
\centering
\caption{Expected detection result.}
\begin{tabular}{@{}lr@{}}
\toprule
\textbf{Quantity} & \textbf{Value} \\
\midrule
$\hat{x}_0$    & 128 \\
$\hat{y}_0$    & 100 \\
$\hat{v}_x$    & 1.50 pixels/frame \\
$\hat{v}_y$    & $-$0.75 pixels/frame \\
$\hat{T}$      & 1881.0 \\
\bottomrule
\end{tabular}
\end{table}

% ──────────────────────────────────────────────────────────────
\section{GPU Parallelization Considerations}
\label{sec:gpu}
% ──────────────────────────────────────────────────────────────

This section describes the key considerations for an efficient CUDA
implementation. Candidates are expected to address \emph{at least} items
\ref{item:thread_map}--\ref{item:reduction}; items
\ref{item:shared}--\ref{item:advanced} are differentiators for strong
candidates.

\subsection{Thread-to-Work Mapping}
\label{item:thread_map}

The search has a natural 4D structure: $(x_0, y_0, v_x, v_y)$. A common
and effective strategy is to assign one GPU thread per starting pixel
$(x_0, y_0)$, with each thread looping over all velocity hypotheses.
This gives $256 \times 256 = 65{,}536$ threads --- sufficient for good
occupancy on modern GPUs.

\begin{notebox}[Thread Block Configuration]
A $16 \times 16$ thread block (256 threads) is a natural choice matching
the 2D spatial structure. With $16 \times 16$ blocks in the grid, this
maps directly to the $256 \times 256$ pixel space.
\end{notebox}

\subsection{Memory Hierarchy}
\label{item:memory}

\begin{itemize}[leftmargin=2em]
    \item \textbf{Constant memory} (\texttt{\_\_constant\_\_}): Ideal for
    the velocity grid. All threads read the \emph{same} velocity values
    simultaneously (broadcast pattern), making constant memory's single-broadcast
    cache architecture a perfect fit. The velocity grid ($25 \times 4$
    bytes $= 100$ bytes) fits trivially in the 64\,KB constant memory limit.

    \item \textbf{Read-only / texture cache} (\texttt{\_\_ldg()}): Frame
    data is read-only and accessed with spatial locality (neighboring threads
    read nearby pixels). Using \texttt{\_\_ldg()} intrinsics routes reads
    through the read-only data cache, which is optimized for this pattern.

    \item \textbf{Registers}: The running sum $T$, validity flag, and
    best-so-far tracking variables should reside in registers. With $N_v^2 = 625$
    iterations per thread, register pressure is low.
\end{itemize}

\subsection{Global Memory Access Patterns}

For adjacent threads (neighboring $x_0$ values), the predicted pixel positions
at frame $t$ are:
\begin{equation}
    p_x^{(i)} = \mathrm{round}(x_0^{(i)} + v_x \cdot t)
\end{equation}

Since $x_0$ values differ by 1 pixel, the read addresses differ by 1 float
--- this is \textbf{coalesced} for small velocities and nearly coalesced for
larger ones, which is favorable for global memory throughput.

\subsection{Reduction}
\label{item:reduction}

After the search kernel, each pixel has a local best statistic. A parallel
reduction is needed to find the global maximum across all $256 \times 256$
pixels. This can be implemented as:
\begin{enumerate}[leftmargin=2em]
    \item A shared-memory tree reduction within each thread block
    \item A final reduction over block-level results (on host or with a
    second kernel)
\end{enumerate}

\subsection{Shared Memory Tiling (Advanced)}
\label{item:shared}

For each velocity hypothesis and frame $t$, threads in a block access a
\emph{shifted} tile of the frame. One could preload the relevant tile
(plus halo) into shared memory before the inner loop. However, given that
the \texttt{\_\_ldg()} cache already handles this pattern well and the frame
data (4\,MB) fits comfortably in L2 cache on modern GPUs, the benefit
may be marginal. A candidate who \emph{analyzes} this tradeoff
(with or without implementing it) demonstrates strong understanding.

\subsection{Further Optimizations (Advanced)}
\label{item:advanced}

\begin{itemize}[leftmargin=2em]
    \item \textbf{Early termination}: Skip velocity hypotheses where the
    first few frames already guarantee the trajectory exits the image.

    \item \textbf{Velocity-parallel decomposition}: Instead of one thread
    per pixel looping over velocities, launch threads over
    $(x_0, y_0, v_x)$ and loop over $v_y$ (or vice versa), increasing
    parallelism at the cost of a more complex reduction.

    \item \textbf{Half-precision accumulation}: Use \texttt{\_\_half2}
    arithmetic for frame reads, accumulating in FP32. This doubles memory
    bandwidth utilization on architectures supporting native FP16.

    \item \textbf{Warp-level primitives}: Use \texttt{\_\_shfl\_down\_sync()}
    for the first stages of reduction, avoiding shared memory bank conflicts.

    \item \textbf{Streams and overlap}: Overlap data transfer with
    computation using CUDA streams (less relevant here since the data is
    small, but shows awareness of the technique).
\end{itemize}

% ──────────────────────────────────────────────────────────────
\section{Performance Targets}
\label{sec:performance}
% ──────────────────────────────────────────────────────────────

\begin{table}[h]
\centering
\caption{Performance reference points.}
\label{tab:perf}
\begin{tabular}{@{}lrrl@{}}
\toprule
\textbf{Implementation} & \textbf{Est.\ Time} & \textbf{Speedup} & \textbf{Notes} \\
\midrule
Python (pure loops)      & $\sim$3 hours     & $1\times$        & Algorithm~\ref{alg:tbd} verbatim \\
NumPy (vectorized)       & $\sim$3 seconds   & $\sim$3600$\times$  & Vectorize over $(x_0, y_0)$ \\
CUDA (basic)             & $<$ 100 ms        & $>$100{,}000$\times$ & Correct thread mapping \\
CUDA (optimized)         & $<$ 20 ms         & $>$500{,}000$\times$ & Memory hierarchy tuned \\
\bottomrule
\end{tabular}
\end{table}

\begin{warnbox}[Correctness First]
An implementation that produces the correct detection in 200\,ms is
\textbf{far more valuable} than one that runs in 5\,ms but gives the wrong
answer. Verify correctness against \texttt{tbd\_reference.txt} before
optimizing.
\end{warnbox}

% ──────────────────────────────────────────────────────────────
\section{Evaluation Rubric}
\label{sec:rubric}
% ──────────────────────────────────────────────────────────────

\begin{table}[h]
\centering
\caption{Scoring rubric (100 points total).}
\label{tab:rubric}
\begin{tabular}{@{}p{4.5cm}cp{7.5cm}@{}}
\toprule
\textbf{Category} & \textbf{Points} & \textbf{Criteria} \\
\midrule
Correctness &
30 &
Detected trajectory matches reference. Proper boundary handling.
No race conditions or undefined behavior. \\
\addlinespace
Performance &
25 &
Kernel time under 100\,ms (baseline pass). Bonus for sub-20\,ms.
Proper use of CUDA events for timing. \\
\addlinespace
Parallelization Strategy &
20 &
Sensible thread-to-work mapping with good occupancy. Block size
justified. Scalability to larger problems. \\
\addlinespace
Memory Optimization &
15 &
Use of constant memory for velocity grid. Awareness of coalescing.
Use of \texttt{\_\_ldg()} or texture cache. \\
\addlinespace
Code Quality &
10 &
Error handling (\texttt{cudaGetLastError}). Readable code structure.
Clean build with no warnings. Comments on design choices. \\
\bottomrule
\end{tabular}
\end{table}

% ──────────────────────────────────────────────────────────────
\section{Context: Why This Matters}
\label{sec:context}
% ──────────────────────────────────────────────────────────────

This interview problem is a simplified version of the real computational
challenge at Zetascope. In the operational system:

\begin{itemize}[leftmargin=2em]
    \item Images are $4096 \times 4096$ (or larger) from qCMOS sensors
    \item Velocity grids may contain $\sim$10{,}000+ hypotheses
    \item Processing must happen in \emph{real-time} within the sensor's
          integration window
    \item The system runs on space-qualified embedded GPUs
          (NVIDIA Jetson AGX Orin class) with strict power and thermal budgets
    \item Objects of interest may contribute fewer than \textbf{one photon
          per frame} on average, requiring integration across hundreds of frames
\end{itemize}

\noindent Efficiently mapping brute-force TBD to GPU hardware is the
foundation upon which all of our more advanced algorithmic accelerations
(temporal-velocity pyramids, multi-resolution search, sensor-algorithm
co-optimization) are built. The skills tested here --- thread mapping,
memory hierarchy exploitation, parallel reduction, and performance
measurement --- are directly applicable to the production system.

\vfill
\begin{center}
\rule{0.5\textwidth}{0.4pt}\\[0.5em]
{\small\textcolor{gray}{%
    Zetascope, Inc.\ --- Confidential Interview Material\\
    Do not distribute.
}}
\end{center}

\end{document}
